% ─────────────────────────────────────────────────────────────────────────────
% ตั้งค่าเอกสารภาษาไทย สำหรับ XeLaTeX
%
% ทำไมต้อง XeLaTeX ไม่ใช่ pdfLaTeX:
%   pdfLaTeX จัดการ Unicode และฟอนต์ระบบไม่ได้ ภาษาไทยจึงต้องใช้ XeLaTeX หรือ LuaLaTeX
%
% ทำไมต้อง \XeTeXlinebreaklocale:
%   ภาษาไทยไม่เว้นวรรคระหว่างคำ ถ้าไม่บอก TeX ว่าเป็นภาษาไทย
%   มันจะไม่ยอมตัดบรรทัดเลย ทำให้ข้อความล้นออกนอกหน้ากระดาษ
% ─────────────────────────────────────────────────────────────────────────────

\usepackage{fontspec}
\usepackage[a4paper,margin=2.4cm,bottom=2.6cm]{geometry}
\usepackage{graphicx}
\usepackage{booktabs}
\usepackage{xcolor}
\usepackage{enumitem}
\usepackage{titlesec}
\usepackage{fancyhdr}
\usepackage{caption}
\usepackage[most]{tcolorbox}
\usepackage{amsmath}
\usepackage[hidelinks]{hyperref}

% ----------------------------------------------------
%  ฟอนต์
%
%  ใช้ฟอนต์ตระกูล TLWG ซึ่งมีทั้งอักขระไทยและละตินอยู่ในไฟล์เดียวกัน
%  จงใจเลือกแบบนี้แทนการจับคู่ Noto Serif (ละติน) กับ Noto Serif Thai
%  เพราะ Noto Serif Thai ไม่มีอักขระละตินเลย การผสมสองฟอนต์จึงต้องพึ่ง
%  ucharclasses ซึ่งสลับฟอนต์ตาม "บล็อกยูนิโคด" ที่เจอ ปัญหาคือมันไม่รู้จัก
%  ขอบเขตของกลุ่ม TeX ดังนั้นข้อความละตินที่ตามหลังกลุ่มอย่าง \textbf{...}
%  จะยังติดฟอนต์ไทยและหายไปทั้งหมด ทั้งยังชนกับ titlesec จนคอมไพล์ไม่ผ่าน
%  ฟอนต์สองสคริปต์ในไฟล์เดียวตัดปัญหานี้ทิ้งทั้งชุด
%
%    Laksaman — ไทยไม่มีหัว + ละตินตัวไม่มีเชิง อ่านง่ายที่สุดสำหรับเนื้อความ
%    Garuda   — ตัวไม่มีเชิงที่หนากว่า ใช้กับหัวข้อและป้ายในไดอะแกรม
%    Tlwg Mono — ความกว้างคงที่จริง ใช้กับชื่อไฟล์และโค้ด
% ----------------------------------------------------

\setmainfont{Laksaman}
\setsansfont{Garuda}
\newfontfamily\thaimono{Tlwg Mono}

\XeTeXlinebreaklocale "th"
\XeTeXlinebreakskip = 0pt plus 0.15pt

% ─── สี ───
\definecolor{ink}{HTML}{1A1A1A}
\definecolor{accent}{HTML}{C0392B}
\definecolor{blue}{HTML}{2C5AA0}
\definecolor{soft}{HTML}{F4F6F8}
\definecolor{warnbg}{HTML}{FDF3F2}
\definecolor{keybg}{HTML}{EEF3FA}
\definecolor{grey}{HTML}{6B7280}

\color{ink}

% ─── หัวข้อ ───
\titleformat{\section}{\sffamily\Large\bfseries\color{blue}}{\thesection}{0.6em}{}
\titleformat{\subsection}{\sffamily\large\bfseries}{\thesubsection}{0.6em}{}
\titlespacing*{\section}{0pt}{2.2ex plus 1ex minus .2ex}{1.2ex plus .2ex}

\setlist[itemize]{leftmargin=1.3em,itemsep=0.25ex,topsep=0.5ex}
\setlist[enumerate]{leftmargin=1.5em,itemsep=0.25ex,topsep=0.5ex}

\captionsetup{font=small,labelfont={bf,color=blue},skip=4pt}

\pagestyle{fancy}
\fancyhf{}
\fancyhead[L]{\small\sffamily\color{grey}\leftmark}
\fancyfoot[C]{\small\thepage}
\renewcommand{\headrulewidth}{0.3pt}

% ─── กล่องนิยามศัพท์ ───
%
% toptitle เผื่อที่ไว้ด้านบนของแถบหัวกล่อง เพราะ tcolorbox คิดความสูงจาก
% ระยะขึ้น-ลงของอักษรละติน ซึ่งไม่เผื่อวรรณยุกต์และสระบนของไทยที่ซ้อนสูงกว่า
% ถ้าไม่ใส่ ไม้โท/ไม้ไต่คู้ในหัวข้อจะถูกตัดบน
\newtcolorbox{term}[1]{
  colback=keybg, colframe=blue, boxrule=0.7pt, arc=2pt,
  left=8pt, right=8pt, top=6pt, bottom=6pt,
  toptitle=3pt, bottomtitle=1pt,
  title={\sffamily\bfseries #1}, coltitle=white, fonttitle=\small,
  breakable
}

% ─── กล่องคำเตือน / กับดัก ───
\newtcolorbox{trap}[1]{
  colback=warnbg, colframe=accent, boxrule=0.7pt, arc=2pt,
  left=8pt, right=8pt, top=6pt, bottom=6pt,
  toptitle=3pt, bottomtitle=1pt,
  title={\sffamily\bfseries #1}, coltitle=white, fonttitle=\small,
  breakable
}

% ─── กล่องสรุปประเด็นสำคัญ ───
\newtcolorbox{keypoint}{
  colback=soft, colframe=grey!45, boxrule=0.5pt, arc=2pt,
  left=8pt, right=8pt, top=6pt, bottom=6pt,
  breakable
}

% ─── ตัวช่วยเล็ก ๆ ───
\newcommand{\eng}[1]{\textsf{#1}}
\newcommand{\code}[1]{{\thaimono\small #1}}
\newcommand{\figref}[1]{รูปที่~\ref{#1}}

% เปลี่ยนคำที่ LaTeX ใส่ให้เองเป็นภาษาไทย
% (ปกติ polyglossia จัดการให้ แต่ที่นี่ใช้แค่ไม่กี่คำจึงเปลี่ยนตรง ๆ ง่ายกว่า)
\renewcommand{\contentsname}{สารบัญ}
\renewcommand{\figurename}{รูปที่}
\renewcommand{\tablename}{ตารางที่}
